\documentclass[11pt]{article}
\usepackage[a4paper,margin=25mm]{geometry}
\usepackage{amsmath,amssymb,amsthm,url}
\newtheorem{theorem}{Theorem}
\newtheorem{lemma}{Lemma}
\newcommand{\R}{\mathbb R}
\newcommand{\dd}{\mathrm d}
\newcommand{\atan}{\operatorname{arctan}}
\newcommand{\cothh}{\operatorname{coth}}
\title{Two nested power-mean inequalities for a logarithm--arctangent mean}
\author{Research manuscript}
\date{29 October 2026}
\begin{document}
\maketitle
\begin{abstract}
We give an exact computer-assisted proof of the two displayed inequalities in Mihokovi\'c--Ra\"issouli's published Conjecture 37. A classical Stolarsky stabilization identity reduces the problem to strict increase of a centered logarithmic ratio. The new comparison is established by a finite, exact 24th-derivative certificate, Abel summation and ordinary calculus. All finite arithmetic is reproducible over $\mathbb Q(\sqrt {17})$; no sampled inequality is used. A separate Lean certificate proves the literal two-branch statement without residual hypotheses. Independent mathematical and formal review is complete; no worldwide priority claim is made.
\end{abstract}
\section{Statement and attribution}
For positive $a,b$, put $M(a,a)=a$, and, when $a\ne b$, put
\[
 M(a,b)=\frac{b-a}{\sqrt2\,\atan((\log b-\log a)/\sqrt2)}.
\]
For $r\ne0$ define $B_r(a,b)=((a^r+b^r)/2)^{1/r}$. Set
\[
 d=\sqrt {17},\qquad \alpha=(5+d)/2,\qquad \beta=(5-d)/2.
\]
Thus $\alpha+\beta=5$, $\alpha-\beta=d$, $\alpha\beta=2$ and $0<\beta<\alpha$.
\begin{theorem}\label{main}
For every $a,b>0$, $a\ne b$, both inequalities hold:
\begin{align*}
M(a,b)&<B_d\bigl(M(a,B_\beta(a,b)),M(B_\beta(a,b),b)\bigr),\\
M(a,b)&<B_{-d}\bigl(M(a,B_\alpha(a,b)),M(B_\alpha(a,b),b)\bigr).
\end{align*}
\end{theorem}
This is the literal displayed claim of \cite{MR}, Conjecture 37, printed p.~21; equation (5), p.~4, defines $M=M_2$. The resultant nesting is defined on p.~1. The earlier arXiv version labels the question Conjecture 4.7. Equality holds on the diagonal. Definition 1, p.~2, separately requires a between-means condition for its terminology ``super-stabilizable.'' That extra condition is outside this manuscript's formal claim. The two branches constitute one outcome.

The Stolarsky identity, finite Abel summation, power-mean monotonicity and calculus used below are classical tools and receive no separate discovery credit. The supplied prior-body reviews exclude several concrete proposed subsumers, but do not certify worldwide novelty.
\newpage
\section{The classical stabilization and ratio transfer}
For $a<b$, define the positive comparison mean
\[
S(a,b)=\left[\frac\beta\alpha\frac{b^\alpha-a^\alpha}{b^\beta-a^\beta}\right]^{1/d}.
\]
Extend it symmetrically, with diagonal value $a$. It is homogeneous. Its stabilization is classical Stolarsky stabilization \cite{R}; the needed identity has the following direct proof. Let $(p,q)=(d,\beta)$ or $(-d,\alpha)$, and let $r=p+q$. Then $\{r,q\}=\{\alpha,\beta\}$, and
\[
 S(a,b)^p=\frac qr\frac{b^r-a^r}{b^q-a^q}.
\]
Set $c=B_q(a,b)$. Since $q>0$, $a<c<b$, and
\[
 c^q-a^q=b^q-c^q=(b^q-a^q)/2.
\]
Adding the corresponding formulas on $(a,c)$ and $(c,b)$ gives
\[
 \frac{S(a,c)^p+S(c,b)^p}{2}=S(a,b)^p,
 \qquad B_p(S(a,c),S(c,b))=S(a,b).                 \tag{1}
\]
Injectivity of positive real powers proves the second equality also for $p<0$.

For $t>0$ introduce centered logarithms
\begin{align*}
 f(t)&=\frac1d\log\frac{\beta\sinh(\alpha t)}{\alpha\sinh(\beta t)},\\
 g(t)&=\log\frac{\sqrt2\sinh t}{\atan(\sqrt2t)},\qquad k(t)=f(t)-g(t).
\end{align*}
These are respectively $\log S(e^{-t},e^t)$ and $\log M(e^{-t},e^t)$; their continuous limits at zero are zero.

We will prove $k'(t)>0$ on $(0,\infty)$. To see that this suffices, let
\[
 t=\tfrac12\log(b/a),\quad x=\tfrac12\log(c/a),\quad y=\tfrac12\log(b/c).
\]
Then $0<x,y<t$ and $x+y=t$. Symmetry and homogeneity give
\[
 M(a,c)=S(a,c)e^{-k(x)}>S(a,c)e^{-k(t)},
 \quad M(c,b)>S(c,b)e^{-k(t)}.
\]
The power mean is homogeneous and strictly increasing in each positive variable for every $p\ne0$. Indeed its derivative in $u$ is
\[
 \frac12u^{p-1}B_p(u,v)^{1-p}>0.
\]
Apply it to the two comparisons and use (1). The resulting lower bound is $e^{-k(t)}S(a,b)=M(a,b)$, proving both inequalities. Symmetry handles $b<a$. Mere order $M<S$ would not justify this transfer; strict increase of the ratio is essential.
\newpage
\section{A differential comparison without inverse tangent}
Define, for real $t$,
\begin{align*}
 A(t)&=\tfrac12\bigl(\cosh(4t)-\cosh t\cosh(dt)+(5/d)\sinh t\sinh(dt)\bigr),\\
 B(t)&=\tfrac12\sinh t\bigl(\cosh(5t)-\cosh(dt)\bigr),\\
 U&=B'A-BA'-A^2,\qquad V=-4BA,\qquad E=(1+2t^2)U+tV.
\end{align*}
Hyperbolic addition formulas and $\alpha\beta=2$ give
\[
 B=\sinh t\sinh(\alpha t)\sinh(\beta t)>0,\qquad
 \frac AB=\cothh t-\frac{\alpha\cothh(\alpha t)-\beta\cothh(\beta t)}d
 =\cothh t-f'(t).                                      \tag{2}
\]
All positivity assertions here and below concern $t>0$. To prove $A>0$ without a series sign assumption, put $a_0=(5/d-1)/4>0$. Then
\[
 A=a_0\bigl(\cosh((d+1)t)-\cosh((d-1)t)\bigr)
 +\tfrac12\bigl(\cosh(4t)-\cosh((d-1)t)\bigr)>0,
\]
since $4<d<5$ and $\cosh$ increases strictly on positive arguments. Expansion at zero gives $A/t^2\to2$ and $B/t^3\to2$. These limits also follow by dividing the identity
\[
 A=\frac{5-d}{2d}\sinh(dt)\sinh t
   +\sinh\bigl((d+3)t/2\bigr)\sinh(\beta t)
\]
by $t^2$ and using $\sinh(ct)/t\to c$.

Set $T=1+2t^2$ and $J=\sqrt2 B/(TA)$. Thus $J\to0$ as $t\downarrow0$. Exact quotient differentiation gives
\[
 J'(t)-\frac{\sqrt2}{T}
   =\frac{\sqrt2 E(t)}{T^2A(t)^2}.                       \tag{3}
\]
Sections 4--5 establish $E(t)>0$ globally. Hence $H=J-\atan(\sqrt2t)$ has positive derivative and zero right limit at zero. It follows that $H(t)>0$: monotonicity gives $H(t/2)\ge0$ from the limit, and $H(t)>H(t/2)$. Since all quantities divided by are positive,
\[
 J>\atan(\sqrt2t)
 \quad\Longrightarrow\quad
 k'(t)=\frac{\sqrt2}{T\atan(\sqrt2t)}-\frac AB>0.       \tag{4}
\]
This supplies exactly the ratio hypothesis used in Section 2. In particular no positivity residual or unproved monotonicity premise is part of Theorem~\ref{main}.
\newpage
\section{Exact finite certificate: data and reconstruction}
We specify a compact exact rule for all certificate data. It also permits independent reproduction without Lean. All operations below take place in $K=\mathbb Q[d]/(d^2-17)$, interpreted at the positive real root. Write $e_\lambda(t)=e^{\lambda t}$, and use the rules
\[
 e_\lambda e_\mu=e_{\lambda+\mu},\quad
 D(t^je_\lambda)=j t^{j-1}e_\lambda+\lambda t^je_\lambda,
 \quad C_\lambda=(e_\lambda+e_{-\lambda})/2,\quad
 H_\lambda=(e_\lambda-e_{-\lambda})/2.
\]
Construct the finite maps
\[
\begin{gathered}
 A=(C_4-C_1C_d+(5/d)H_1H_d)/2,\quad B=H_1(C_5-C_d)/2,\\
 U=(DB)A-B(DA)-A^2,\qquad V=-4BA.
\end{gathered}
\]
Combine equal exponents by addition. Thus the rule completely determines each coefficient $u_\lambda,v_\lambda$ in $U,V$. The constant coefficients are $u_0=-87/68$, $v_0=0$; $u_{-\lambda}=u_\lambda$, $v_{-\lambda}=-v_\lambda$. The positive frequencies, in strictly decreasing order, are
\[
\begin{gathered}
7+d,\ 2+2d,\ 10,\ 5+d,\ 2d,\ 8,\\
3+d,\ -2+2d,\ 7-d,\ 2,\ -3+d,\ 5-d.
\end{gathered}                                                     \tag{5}
\]
There are no other nonzero frequencies. For each frequency in (5), set
\begin{align*}
 c_{i0}&=\lambda_i^{22}(u_i\lambda_i^2+24v_i\lambda_i+1104u_i),\\
 c_{i1}&=\lambda_i^{22}(v_i\lambda_i+94u_i),\qquad
 c_{i2}=2\lambda_i^{22}u_i.
\end{align*}
For $j=0,1,2$ and $1\le m\le12$ form $P_{mj}=\sum_{i=1}^m c_{ij}$. All 36 prefixes are strictly positive. A reproducible rational certificate is: reduce each $P_{mj}=a+bd$ to a rational pair and check
\[
 a+b\ell_b>0,\qquad
 \ell_b=\begin{cases}41/10&b\ge0,\\21/5&b<0.\end{cases}             \tag{6}
\]
The bounds follow from $(41/10)^2<17<(21/5)^2$ and also verify the order (5). The last three sums are
\[
\begin{array}{r}
355707128874106034847744,\\
46479400985961399582720,\\
1463704104369024663552.
\end{array}
\]
The supplement contains every pair and lower bound in \texttt{exact-results.json}, generated by \texttt{exact\_check.py}. That standard-library checker implements precisely the finite rules above, checks support, signs, all prefixes, the derivative identity and jets, and rejects a deliberate 94-for-96 derivative mutation. No floating point is used.
\newpage
\section{Why the finite certificate is a global proof}
Repeated Leibniz differentiation, pairing positive and negative frequencies, yields the exact identity
\[
 \frac{E^{(24)}(t)}2=\sum_{i=1}^{12}\left[
 c_{i0}\cosh(\lambda_it)+t(c_{i1}+c_{i2})\lambda_i\sinh(\lambda_it)
 +t^2c_{i2}\lambda_i^2\cosh(\lambda_it)\right].                \tag{7}
\]
For verification, the coefficients of $t^0,t^1,t^2$ at a positive frequency before pairing are respectively
\[
 \lambda^{24}u+24\lambda^{23}v+1104\lambda^{22}u,
 \quad\lambda^{24}v+96\lambda^{23}u,\quad2\lambda^{24}u.
\]
In particular the linear term uses $c_1+c_2$, not $c_1$ alone. The constant contribution $-(87/68)(1+2t^2)$ is annihilated only after taking 24 derivatives.

For any finite array, let $S_m=\sum_{i=1}^m a_i$. Abel's identity is
\[
 \sum_{i=1}^n a_iw_i=S_nw_n+\sum_{i=1}^{n-1}S_i(w_i-w_{i+1}).
\]
If all prefixes are positive and weights descend nonnegatively, the sum is nonnegative; it is positive if $w_n>0$. For $t\ge0$, each of
\[
 \cosh(\lambda t),\qquad\lambda\sinh(\lambda t),\qquad
 \lambda^2\cosh(\lambda t)
\]
is increasing and nonnegative in $\lambda>0$. Their derivatives in $\lambda$ are, respectively,
\[
 t\sinh(\lambda t),\quad\sinh(\lambda t)+\lambda t\cosh(\lambda t),
 \quad2\lambda\cosh(\lambda t)+\lambda^2t\sinh(\lambda t).
\]
Apply Abel to $c_0$, $c_1+c_2$ and $c_2$ using (5)--(6). The first weighted sum is strictly positive; the others are nonnegative. Thus (7) proves $E^{(24)}(t)>0$ for every $t\ge0$.

The jets $E^{(j)}(0)$ for $0\le j<24$ vanish except for
\[
\begin{array}{c|r}
j&E^{(j)}(0)\\\hline
10&37847040\\12&19891322880\\14&5550283554816\\
16&1162696473968640\\18&206947471161556992\\
20&33226670114812723200\\22&4978839848478924865536
\end{array}
\]
They follow either from the same finite maps or the integer rules on the next page. Since $E$ is smooth, for $j=23,22,\ldots,0$ the identity
\[
 E^{(j)}(t)=E^{(j)}(0)+\int_0^t E^{(j+1)}(s)\,\dd s
\]
proves $E^{(j)}(t)>0$ for every $t>0$, by descending induction. This proves the required global positivity, then (3)--(4), and finally Theorem~\ref{main}.
\newpage
\section{Integer cross-check and literal formal scope}
For completeness, a second rule independently checks the jets. Write $A_n=A^{(n)}(0)$, $B_n=B^{(n)}(0)$. For even $n$,
\[
2A_n=4^n-\sum_{\substack{0\le j\le n\\j\text{ even}}}\binom nj17^{j/2}
 +5\sum_{\substack{1\le j\le n\\j\text{ odd}}}\binom nj17^{(j-1)/2},\qquad B_n=0.
\]
For odd $n$, $A_n=0$ and
\[
 2B_n=\sum_{\substack{0\le j\le n\\j\text{ even}}}\binom nj(5^j-17^{j/2}).
\]
With negative indices zero, let
\begin{align*}
 U_n&=\sum_{j=0}^n\binom nj
 (B_{j+1}A_{n-j}-B_jA_{n-j+1}-A_jA_{n-j}),\\
 V_n&=-4\sum_{j=0}^n\binom njB_jA_{n-j}.
\end{align*}
Then $E^{(n)}(0)=U_n+nV_{n-1}+2n(n-1)U_{n-2}$. The supplied checker cross-checks these values through order 24 against direct exponential differentiation. This is exact integer/rational arithmetic, not evidence from a finite numerical grid.

The immutable Lean mathematical sources comprise 35 modules. Their final declaration is
\begin{center}
\texttt{LogArctangentSource.source\_conjecture37 : SourceClaim}.
\end{center}
The actual \texttt{SourceClaim} universally quantifies real $a,b$, assumes only $0<a$, $0<b$, $a\ne b$, and concludes the conjunction of the two inequalities in Theorem~\ref{main}. The mean is the literal off-diagonal quotient above, with an if-branch assigning the diagonal value. The power mean uses real powers, with a separate geometric branch at exponent zero; neither parameter here is zero. No residual $E$-positivity or monotonicity hypothesis appears in the endpoint.

The packet verifier records 561 module-data entries covering 560 unique kernel declarations, including generated, private and instance declarations. One generated equation is listed twice; its actual origin is checked separately. It records actual types, transitive axioms and safety flags. The only permitted transitive axioms are \texttt{propext}, \texttt{Classical.choice}, \texttt{Quot.sound}. Generated partial recursors are reported by their actual flags, without treating a name containing ``unsafe'' as an unsafe declaration.

The ordinary proof above is independent of interpreting a successful replay as a mathematical exposition. Conversely, replay proves no source priority or novelty claim. The optional ordinary between-means argument retained in the source materials is not formalized in this endpoint and is not asserted here.
\newpage
\section{Review and reproducibility limits}
Independent final review and source-only replay passed for the displayed two-branch theorem. A current public-registry audit completed on 29 October 2026 at 08:10 UTC found no matching prior proof within its documented coverage. The most substantial remaining objection is potential subsumption by a prior same-object global ratio theorem or a verified classical specialization. Missing MathDB access, incomplete specialist-body coverage and bounded registry screening remain limits; no worldwide priority is certified.

The source-only archive contains no oleans. Replay uses Lean 4.34.1 and eight pinned packages, serial compilation, warning-as-error and explicit binders. Default-parallel compilation previously reached the 3 GiB limit in the root run; serial compilation is required, without raising that limit. Packet documentation records measured memory, wall deadlines, binary hash, source hashes and extracted replay results. These controls are evidence for the stated bytes, not acceptance authority.
\begin{thebibliography}{9}
\bibitem{MR} L. Mihokovi\'c and M. Ra\"issouli,
\emph{Results on comparison and sub/super-stabilizability of some new means},
Applicable Analysis and Discrete Mathematics \textbf{20}(1) (2026), 1--23.
DOI: \url{https://doi.org/10.2298/AADM240618001M}.
Published Conjecture 37, p.~21; earlier arXiv:2405.13618v2, Conjecture 4.7.
\bibitem{R} M. Ra\"issouli, \emph{Stabilizability of the Stolarsky mean and its approximation in terms of the power binomial mean}, manuscript received 25 May 2011, Theorem 3.1.\newline\url{https://rgmia.org/papers/v14/v14a43.pdf}.
\end{thebibliography}
\end{document}
